% !TEX root = task-2-7-2.tex
% Editable task bank. Append new BankTask blocks at the end of the document.
\documentclass[12pt,a4paper]{article}
\usepackage{physics-textbook}
\usepackage{graphicx,needspace}
\usetikzlibrary{decorations.pathmorphing}
\newif\ifShowSolutions
\ifdefined\TutorEdition
  \ShowSolutionstrue
\else
  \ShowSolutionsfalse % Change to \ShowSolutionstrue to display solutions.
\fi
\newcounter{banktask}
\newsavebox{\BankTaskBox}
\newcommand{\BankTaskContent}[3]{%
  \noindent\textbf{\thebanktask.}\quad #2\par
  \ifShowSolutions
    \Needspace{50mm}% Keep the solution heading with its opening text or diagram.
    #3\par
    \ifdefined\TutorEdition
      \ifcsname TutorSource#1\endcsname
        \nopagebreak\smallskip{\small\csname TutorSource#1\endcsname\par}%
      \fi
    \fi
  \fi
}
\newcommand{\BankTask}[3]{%
  \refstepcounter{banktask}\par\medskip
  \sbox{\BankTaskBox}{\begin{minipage}{\linewidth}%
    \BankTaskContent{#1}{#2}{#3}\end{minipage}}%
  \ifdim\dimexpr\ht\BankTaskBox+\dp\BankTaskBox\relax>\dimexpr\textheight-20mm\relax
    \Needspace{35mm}\BankTaskContent{#1}{#2}{#3}%
  \else
    \Needspace{\dimexpr\ht\BankTaskBox+\dp\BankTaskBox\relax}%
    \noindent\usebox{\BankTaskBox}\par
  \fi
  \medskip
}
\begin{document}
\BookSubsection{2.7.2}{Динамическое равновесие}


% Source 4DBED7: https://ege.fipi.ru/bank/questions.php?search=1&proj=BA1F39653304A5B041B656915DC36B38&qid=4DBED7
\BankTask{4DBED7}{
\begin{center}
\begin{tikzpicture}[>=stealth,x=0.85cm,y=0.85cm]
\draw[->] (0,0)--(4.5,0) node[below] {$t$};
\draw[->] (0,0)--(0,2.8) node[left] {$n$};
\node[below left] at (0,0) {$0$};
\draw[thick] (0,0.48) .. controls (0.65,1.2) and (1.15,1.85) .. (1.9,1.85)--(4.05,1.85);
\node at (0.65,1.7) {$1$};
\node at (2.8,2.23) {$2$};
\end{tikzpicture}
\end{center}
В стеклянную колбу налили немного воды и герметично закрыли её пробкой. Вода постепенно испарялась. На рисунке показан график изменения со временем \(t\) концентрации \(n\) молекул водяного пара внутри колбы. Температура в колбе в течение всего времени проведения опыта оставалась постоянной.

В конце опыта в колбе ещё оставалась вода.
\begin{enumerate}
\item На участке \(1\) плотность водяных паров уменьшалась.
\item На участке \(2\) давление водяных паров не менялось.
\item На участке \(1\) водяной пар ненасыщенный, а на участке \(2\) насыщенный.
\item На участке \(1\) давление водяных паров уменьшалось.
\item На участке \(2\) плотность водяных паров увеличивалась.
\end{enumerate}
}{
\smallskip\noindent\textbf{Решение от GPT}\par
При постоянной температуре \(p=nkT\), а плотность пара \(\rho=m_0n\), где \(m_0\) — масса молекулы воды.
\begin{enumerate}
\item Неверно: на участке \(1\) концентрация и плотность растут.
\item Верно: на участке \(2\) концентрация постоянна, поэтому давление не меняется.
\item Верно: на участке \(1\) пар ещё ненасыщенный; на участке \(2\) устанавливается равновесие с оставшейся водой, пар насыщенный.
\item Неверно: на участке \(1\) давление растёт вместе с концентрацией.
\item Неверно: на участке \(2\) плотность постоянна.
\end{enumerate}
\(\boxed{23}\).
}

\end{document}
